%-*-latex-*-
%-*-latex-*-
\newcommand\COURSE{ciss240}
\newcommand\ASSESSMENT{q1901}
\newcommand\ASSESSMENTTYPE{Quiz}
\newcommand\POINTS{\textwhite{xxx/xxx}}

\input{myquizpreamble}
\input{yliow}
\input{\COURSE}
%-*-latex-*-

%-*-latex-*-

%-*-latex-*-
\renewcommand\TITLE{\ASSESSMENTTYPE \ \ASSESSMENT}

\renewcommand\EMAIL{}
\renewcommand\AUTHOR{eagraham1@cougars.ccis.edu}

\textwidth=6in
\begin{document}
\topmatterthree

This is a closed-book, no compiler, 5 minute quiz.

%------------------------------------------------------------------------------
\nextq
Complete the following function to check if an integer is prime.
For instance \verb!isprime(2)! returns true,
\verb!isprime(-3)! returns false,
\verb!isprime(1)! returns false,
and
\verb!isprime(10)! return false,
\\
\ANSWER
\begin{answercode}
// isprime(n) return true exactly when n is a prime.
bool isprime(int n)
{
    if (n <= 1) return false;
    for (int i = 2; i * i <= n; ++i)
    {
        if (n % i == 0) return false;
    }
    return true;
}

\end{answercode}

%------------------------------------------------------------------------------
\nextq
Complete the following function to compute powers.
For instance \verb!mypow(2, 0)! returns 1 and
\verb!mypow(2, 4)! returns 16.
\\
\ANSWER
\begin{answercode}
// mypow(x, y) return x to the power of y. Both x and y are integers.
// Assume y >= 0. Obviously you are not allowed to use the pow function.
// HINT: You need a loop.
int mypow(int x, int y)
{
    int result = 1;
    for (int i = 0; i < y; ++i)
    {
        result *= x;
    }
    return result;
}
\end{answercode}

%------------------------------------------------------------------------------
\nextq
Complete the following function that returns the number of digits in an integer.
For instance
\verb!numdigits(0)! is 1,
\verb!numdigits(4030201)! is 7,
\verb!numdigits(-4030201)! is 7.
\\
\ANSWER
\begin{answercode}
int numdigits(int n)
{
    if (n < 0) n = -n;
    if (n == 0) return 1;
    int count = 0;
    while (n > 0)
    {
        n /= 10;
        count++;
    }
    return count;
}
\end{answercode}

%------------------------------------------------------------------------------
%-*-latex-*-
\newpage
\input{instructions}
\end{document}

